% !TEX program = xelatex
% 2026年11月绍兴市高考科目诊断性考试 语文试题
\documentclass[11pt,a4paper]{ctexart}
\usepackage{ywexam}

\begin{document}

\examtitle{2026年11月绍兴市高考科目诊断性考试}{语文试题}
\examnote{说明：本试卷的所有答案一律写在“答题纸”上。}

\sechead{一、阅读（73分）}

\subsechead{（一）阅读Ⅰ（19分）}
阅读下面的文字，完成1～5题。

\material{材料一：}
\begin{materialtext}
各位朋友，危机能让人看清真相。而在这里，在伦敦——狄更斯笔下的城市，我们清楚地看到，当今世界正面临“双重危机”：气候危机不断推高气温，使我们日益逼近灾难性的临界点；能源危机暴露了依赖碳氢化合物的世界所犯下的愚蠢错误。表面上看，这两个危机似乎彼此独立，但驱动它们的是同一种破坏性力量——化石燃料。要应对它们，也必须采取同一种方案：快速、公平地向清洁能源转型，同时大力加强适应能力、韧性建设和对已遭受气候危害群体的气候公正。

乱象正在我们眼前加速演进，我们刚刚经历了有记录以来最热的十一年。气候灾害越来越频繁，越来越具破坏性，造成的损失也越来越大。世界气象组织已发出警告，更严峻的情况还在后头。厄尔尼诺不只是在敲门，它可能把整个房子掀翻。气温不断升高，扰乱粮食和水系统，而弱势群体遭受的打击最为沉重。十年前，世界各国领导人在巴黎商定，要将全球气温升幅限制在1.5摄氏度以内。如今，科学家们警告说，未来几年全球年均升温将超过这一阈值。

摆在我们面前的任务是：严格限制温升过冲、缩短过冲持续时间，并尽快将温升拉回1.5摄氏度以下。哪怕是零点几度都至关重要，每一刻都十分关键。因为过冲幅度越大、持续时间越长，突破地球气候临界点的风险就越大。一旦突破，将带来不可逆转的变化：珊瑚礁系统将走向崩溃；格陵兰和西南极冰盖将加速消融，这将导致海平面持续上升，进而重塑海岸线，迫使数以百万计的人流离失所，并威胁一些岛国的生存；调节天气和降雨的主要海洋环流系统将减弱；亚马逊雨林部分地区将逐渐转变为类似稀树草原的环境。地球气候临界点就像汽车后视镜里的物体，它们比看起来要近得多。

与此同时，我们正面临第二个危机，中东冲突引发了前所未有的能源冲击。国际能源署指出，这场冲击的规模堪比1970年代的石油危机……对许多发展中国家来说，这不仅仅是一场能源危机，还是一场债务冲击、粮食冲击、发展冲击。我还要补充一点：任何和平协议都值得欢迎，也将带来亟需的缓解；但可以肯定的是，危机的影响很可能是长期的。

各位朋友，当前的“双重危机”再次暴露了一种过时发展模式的局限性。这种模式依赖化石燃料，而一场冲突就可能颠覆全球能源供应，一个关键节点受阻就会导致价格飙升。这种模式把自然视为取之不尽的资源，可以不计后果地肆意消耗。这种模式创造了巨大财富，但也加深了不平等、加剧了不安全。在这种模式下，那些对造成危机责任最小的群体却付出了最沉重的代价。当务之急是下定决心，全面落实可持续发展目标。

好消息是，不同于以往任何一次能源危机，我们如今已有一条明确的出路，一条清洁的出路。在世界大部分地区，可再生能源是成本最低、获取速度最快、最具规模化潜力的新增电力来源。2010年以来，太阳能成本下降了近90\%，陆上风电成本下降超过70\%，而电池储能成本下降了95\%。去年，风能和太阳能的发电增量已超过全球全部新增电力需求。太阳能创下了历史上所有电力来源中最大年度增幅纪录。全球新增可再生能源发电量中，90\%以上的成本已低于最廉价的化石燃料替代方案。根据国际可再生能源署的数据，仅2025年一年，现有可再生能源发电能力就为全球经济节省了4800亿美元的化石燃料支出。同时，可再生能源避免的二氧化碳排放量超过了美国、欧盟和日本年排放量的总和。此外，清洁能源吸引的投资已接近化石燃料的两倍。这一势头很大程度上来自化石燃料进口国，它们决心摆脱不稳定、不可预测的能源市场。阳光不会遭禁运，风也不会被封锁。

各位朋友，结论已经很明确：能源独立绝不能建立在对化石燃料的依赖上。可再生能源才是真正能源安全的基石。推动交通、建筑和工业电气化，是减少排放、摆脱对进口化石燃料依赖的最快捷途径之一。一个经济体使用清洁电力越多，就越安全、越有韧性、越具竞争力。
\end{materialtext}
\begin{flushright}（摘译自联合国秘书长古特雷斯2026年6月23日在“伦敦气候行动周”上的讲话）\end{flushright}

\material{材料二：}
\begin{materialtext}
全球气候变化治理无法形成一套符合各方诉求和利益的减排方案，发达国家和发展中国家在气候变化问题上据理力争，使得气候变化问题逐渐转变为覆盖全球的政治问题。大国气候博弈的本质是发展权之争。

第一，经济发展和碳减排间的矛盾难以消除。气候变化风险的核心影响要素是碳排放环境风险，然而新型清洁能源和气候变化减缓技术难以在短时间内取代传统能源，也就无法促使那些依赖高碳经济发展模式的国家进行低碳转型。此外，发展中国家的首要治理目标并非应对气候变化风险，而是消除贫困、改善民生和发展经济等方面。

第二，气候变化风险治理需要科学技术的识别和大量资金的投入。气候有益技术供给不足使得大部分发展中国家无法自主应对碳排放环境风险，技术缺口不断加剧。发达国家掌握着大部分气候有益技术，虽然《巴黎协定》中有“技术开发和转让”条款，但国际社会并未对气候有益技术的公共产品属性达成统一认识，进而造成了知识产权保护制度阻碍气候有益技术向发展中国家转让的局面。在资金方面，全球气候治理尚未形成完善的资金机制且面临供资难题，资金缺口较大。《巴黎协定》中的“大规模资金”愿景并未就资金分摊机制进行建构，致使发达国家在兑现承诺时极易出现推诿和延迟现象。资金缺位使得发展中国家难以提升气候变化风险的应对能力。
\end{materialtext}
\begin{flushright}（摘编自周珂《全球气候变化风险治理》）\end{flushright}

\question{1}{下列对材料相关内容的理解和分析，不正确的一项是（3分）}
\begin{choices}
\item 气候危机与能源危机看似独立，实则同由化石燃料驱动，必须以同一方案加以应对。
\item 当前中东冲突引发的能源冲击规模空前，可能让众多发展中国家陷入多重发展困境。
\item 依赖化石燃料的发展模式有局限性，封锁、禁运等国际事件或导致能源供应不稳定。
\item 为了不让能源供应受外部市场波动的影响，一个国家需要充分利用本国的化石燃料。
\end{choices}

\question{2}{材料一是一篇演讲稿，下列对其写作特点的分析，不恰当的一项是（3分）}
\begin{choices}
\item 开篇点明“双重危机”的共同根源，由此提出相关对策，有利于向听众清晰准确地传达自己的观点。
\item 以近年切身可感的“高温”叙事，拉近与听众的距离，引发共鸣，同时突出了气候变化的严重后果。
\item 借“第二个危机”强调能源独立和能源安全的重要性，针对性强，警告各国摆脱对化石燃料的依赖。
\item “阳光不会遭禁运，风也不会被封锁”一句形象而风趣，描绘了摆脱能源市场控制后的经济优势。
\end{choices}

\question{3}{下列现象最能支持材料二观点的一项是（3分）}
\begin{choices}
\item 历史上碳排放最少的贫穷国家，遭遇极端热浪时人均GDP损失却是富裕地区的4.5倍。
\item 对于广大发展中国家来说，关停煤电意味着工厂停工、医院断电、陷入黑暗与贫困。
\item 印度大米出口占全球40\%，2026年热浪导致减产，推高了全球米价，还引发了移民潮。
\item 某能源进口国近年持续扩建本土光伏与风电，能源自给率和抗风险能力有了稳步提升。
\end{choices}

\question{4}{下列文字是材料一演讲的结语，请根据上文，在画横线处补写合适内容，每处字数在30字左右。（4分）}

让我以开头提到的狄更斯作结。对气候议程而言，这确实是“最好的时代”，也是“最坏的时代”。最坏之处在于，\underline{\hspace{0.6em}A\hspace{0.6em}}。但这也是最好的时代，因为\underline{\hspace{0.6em}B\hspace{0.6em}}。我们拥有巨大的机遇，也肩负重大责任。

\question{5}{根据材料所述的危机与困境，某研究所需要提供相应的治理建议，请填充下列表格。（6分）}

\begin{center}
\begin{tabular}{|c|p{6.2cm}|p{6.2cm}|}
\hline
序号 & 危机与困境 & 对策或建议 \\
\hline
1 & 气候危机不断推高气温，使我们日益逼近灾难性的临界点。 & \xq{1} \\
\hline
2 & 能源危机暴露了过时发展模式的局限性，中东冲突引发了前所未有的能源冲击。 & \xq{2} \\
\hline
3 & 全球气候变化治理无法形成一套符合各方诉求和利益的减排方案。 & \xq{3} \\
\hline
\end{tabular}
\end{center}

\subsechead{（二）阅读Ⅱ（18分）}
阅读下面的文字，完成6～9题。

\begin{center}{\bfseries 露天电影}\par\vspace{0.15em}徐则臣\end{center}

\begin{materialtext}
在乡村，露天电影这些年已经绝迹了，城市里又兴盛起来。

大家开着车到露天电影下，很少看电影，很多人围聚在一起，参差不齐的人头安静或者攒动。世界一下子就热闹了，人与人之间就有了一种温暖而又隐秘的关联，即使相互还是陌路，也让我感动。都坐在小板凳上、椅子上，或者席地而坐、爬到树顶，零散各处的人在这里形成了一种和谐的秩序。我喜欢这种日常的、过日子的场景。

小时候一听来电影，晚饭都吃不踏实，心悬着，到处打听今晚放的是什么片子。最好是枪战的、武打的，侦探的也行，一到城市里的谈情说爱就不喜欢。主角的手拉在一起，像什么样子。太不像话了，老头老太太就骂了，把脸扭过去，说他们把爹娘的脸都丢尽了。年轻人起哄，我们这些小屁孩也不懂，也煽风点火地起哄。外国的片子所有人都不喜欢。\uline{\xq{1}长得怪模怪样不说，张嘴就是“亲爱的”，受不了。}没见过哪家的男人对老婆这样说过，也没见哪家的女人这样对丈夫说过。

在搬到这个楼下有公园的地方之前，我差不多十年没看过露天电影。不记得确切的时间了，大概十年前，从外面的城市回到家里，一个假期下来，总觉得少干点事。某一天走过中心路，停下来，想起现在两层楼的地方过去埋着两根杆子，每个月都有那么一两次，杆子上拴着绷紧的银幕。一块四四方方的白帆布，四个角吊起来。一台机器发出人跑步的声音，很多人就着五彩的光柱子跳上了银幕。\uline{\xq{2}他们在我们斜上四十五度的高处生活，过着与我们不一样的日子。}我们像做梦一样看他们，看他们怎么在另外一个个陌生美好的世界里生老病死。那些世界我们曾经是多么地羡慕，一群屁大的孩子跟着放映机，一个一个地跑遍周围的村庄，来来回回看同一场电影。谁都相信，那世界里不止这些幸福的人，这些好玩的事。可是，还有些什么，不知道，所以要一遍遍接着看。

不知道我们那一帮整天跟着放映机跑的孩子里，有谁发现了另外的人事。我是没有，若干年以后，经见了更多的世事，我常会在某一时刻发一下子呆，觉得此情此景恍惚是经历过了的，又找不出证据，就自恋地夸奖一下自己，说不定这就是电影之外的人事，在多年前还是个孩童和少年的时候，就用想象力提前发现了。我乐于这样褒奖自己，顺便重温一下当年的露天电影。我几乎没有错过一场有机会观看的电影，我说的是在中心路上放映的露天电影。它距我家不远，有人在银幕上咳嗽一声，在家里都能准确地判断出他是男还是女。

但是露天电影不再有了。那天我站在埋杆子的地方深切怀念，远游归来这世界变了样，变就变在露天电影没了。不跟我说一声就没了。多少年里，我从跟在父母后头，到自己搬小板凳坐到电影背面，到搬椅子跻身正面，再到更大了，凳子椅子都不屑搬了，站着看，先是挤在人群里站着，然后是一个人站到一边，有点冷清，内心里却颇感到悲壮的孤独，就是现在小孩都在意的，酷。年龄又大了一些，在外地念书了，真觉得自己在露天电影前得有点酷。我跟电影的距离远。而他们，我的父老乡亲，多少年如一日地，还跟电影上的人像一家子。我就在心里说，都是瞎编的，我知道。好像他们就不知道。

现在，我能看见当年脸上的倨傲是多么可笑，那种刻意的躲避和疏离是现在我所鄙视的。我果真不需要和他们在一起吗？恰恰相反，在内心里，我是多么希望能像没有离开家乡时那样，端着饭碗凑到他们堆里，肆无忌惮地把筷子伸到别人的碗里夹点想吃的菜。可是疏离是不可避免的，疏远也是无能为力的。现在每年短暂的还乡，我成了他们的客人。他们用对远道而来的客人一样的目光看我，说，回来了？难道我不该回来吗？不仅他们，就连那些一起追着放映机跑遍村庄的伙伴，也变得客气了。他们说，回来了？我说回来了。他们急于躲开，我急于逃避。怎么会这样？我们见了面突然不知道把手和眼神往哪里放了。我们都像发现了电影里的那个世界之外的人事，它们突如其来地出现，让我们在各自的眼里看到了陌生。这陌生把一起追逐露天电影时心照不宣的热情紧紧地包裹起来。好了，在故乡，我终于成了一个异乡人。

露天电影就这么悄无声息地没了。我问母亲，她也说不清什么时候放了最后一场。那时候露天电影是多么兴盛，婚丧嫁娶，都要包一两场电影以谢乡邻。一周多了能看八场，遇事的两家同时放映。我还能想起当年哪边武打、枪战激烈就往哪边跑的景象，\uline{\xq{3}很多人和我一样，椅子凳子举在头顶像一群阿拉伯人。}整个晚上凳子几乎都坐在我们头上，两边的电影比赛着精彩，根本没有时间把凳子安稳地放下来坐一会儿。如果没有月亮，路就是黑的，天上有星星，中心路上忙得像赶夜集。这样的夜晚，快乐也马不停蹄。

现在，乡村的夜晚安静了，空无一人似的安静，只看见灯光这里亮，那里亮。都关着门看电视，看影碟。多年以前的愿望终于实现了，我们成功地把电影搬到了自己家里。没有人在大街上溜达，闲也闲在家里，不拿出来给人看。都躲在家里，电影放给谁看。不知道放映员是不是这么想的。夜晚一下子倒退了多少年，安静得狗都不敢叫，到处都是空的，端着饭碗出来找不到人的空。

我不知道露天电影消失的时代好不好。想着让我难受。时间久了，忘了疼，只在想起“露天电影”这个美好的短语时，心里抖一下，再抖一下，就去干别的事了。搬到芙蓉里的第三天晚上，晚饭后经过楼下，看见公园广场上挂起来的银幕，身心都抖起来。

我对晓说：“露天电影。”

晓说：“哪来的露天电影？”

“一定是。”

天黑了，我趴在窗台上盯着那块银幕，直到人群聚拢过来，直到声音响起来，直到文字和人在银幕上走起来，\uline{\xq{4}回过头对晓说：“露天电影。嘿嘿。”}
\end{materialtext}
\begin{flushright}（有删改）\end{flushright}

\question{6}{下列对文本相关内容和艺术特色的分析鉴赏，不正确的一项是（3分）}
\begin{choices}
\item 本文以我的成长视角，写出了与故乡逐渐疏离的过程，真切自然，富有抒情气息。
\item 露天电影消失后的安静与兴盛时的热闹对比，强化了公共生活瓦解后乡村的冷清。
\item 本文虚实结合，虚写我们“都关着门看电视，看影碟”，实写青年进城后夜晚的安静。
\item 全文语言质朴隽永，平静的回忆性叙述中寄寓深情，又有对生活变迁的清醒审视。
\end{choices}

\question{7}{关于文中画横线的语句，下列说法不正确的一项是（3分）}
\begin{choices}
\item 句\xq{1}，以诙谐的笔调表现了当时乡村观众对外来情感表达方式的隔膜。
\item 句\xq{2}，以细节强调电影放映条件简陋，反衬它在匮乏年代的精神慰藉。
\item 句\xq{3}，以有趣的联想写人们转场时的急切，场景热闹混乱却不失和谐。
\item 句\xq{4}，以短句和拟声词“嘿嘿”收束全文，有窃喜也有得意，余味悠长。
\end{choices}

\question{8}{“好了，在故乡，我终于成了一个异乡人”一句意蕴丰富，请结合全文谈谈你的理解。（6分）}

\par\vspace{2.2em}

\question{9}{请你以本文为素材，给农村“露天电影”这一事物撰写一则150字左右的“家乡风物志”条目。“志”的表述方式以叙述和说明为主，要写清楚其来源、特点、价值等。（6分）}

\par\vspace{2.2em}

\subsechead{（三）阅读Ⅲ（21分）}
阅读下面的文言文，完成10～14题。

\material{材料一：}
\begin{materialtext}
\wave{昔者秦王赵正出游天下还至柏人而病病笃。}喟然流涕长太息。召丞相斯而告之曰：“吾自视天命，年五十岁而死。吾行年十四而立，立卅七岁矣。吾当以今岁死，而不知其月日，故出游天下，欲以变气易命，不可欤？今病笃，几死矣。其亟日夜输趋，至甘泉之置，毋须后者。其谨微密之，毋令群臣知病。”

病即大甚，而不能前，故复召丞相斯曰：“吾霸王之寿足矣，不奈吾子之孤弱何。吾闻之：牛马斗，而蚊虻死其下；大臣争，齐民苦。哀怜吾子之孤弱，及吾蒙容之民，死且不忘。其议所立。”丞相臣斯、御史臣去疾昧死顿首言曰：“今道远而诏期窘，臣恐大臣之有谋，请立子胡亥为代后。”王曰：“可。”

王死而胡亥立，即杀其兄扶苏、中尉恬。大赦罪人，而免隶臣高以为郎中令。因夷其宗族，坏其社稷，燔其律令及故世之藏。又欲起属车万乘以抚天下，曰：“且与天下更始。”
\end{materialtext}
\begin{flushright}（节选自《赵正书》\xq{1}，有删改）\end{flushright}

\material{材料二：}
\begin{materialtext}
天下失始皇帝，皆\zd{遽}恐悲哀甚，朕奉遗诏，今宗庙事及箸以\zd{明}至治大功德者具矣，律令当除定者毕矣。\uline{元年与黔首更始，尽为解除故罪。}令皆已下矣，朕将自抚天下。吏、黔首其具行事，毋以徭赋扰黔首，毋以细物苛劾县吏。亟布。
\end{materialtext}
\begin{flushright}（节选自《秦二世元年十月甲午诏书》\xq{2}）\end{flushright}

\material{材料三：}
\begin{materialtext}
至平原津而病。始皇恶言死，群臣莫敢言死事。上病益甚，乃为玺书赐公子扶苏曰：“与丧会咸阳而葬。”书已封，在中车府令赵高行符玺事所，未授使者。七月丙寅，始皇崩于沙丘平台。丞相斯为上崩在外，恐诸公子及天下有变，乃秘之，不发丧。棺载轀凉车中，故幸宦者\zd{参乘}，所至上食。百官奏事如故，宦者辄从轀凉车中可其奏事。独子胡亥、赵高及所幸宦者五六人知上死。赵高故尝教胡亥书及狱律令法事，胡亥私幸之。高乃与公子胡亥、丞相斯阴谋破去始皇所封书赐公子扶苏者，而更诈为丞相斯受始皇遗诏沙丘，立子胡亥为太子。更为书赐公子扶苏、蒙恬，数以罪，其赐死。行从直道至咸阳，发丧。太子胡亥袭位，为二世皇帝。
\end{materialtext}
\begin{flushright}（节选自《史记·秦始皇本纪》，有删改）\end{flushright}

\material{材料四：}
\begin{materialtext}
大抵偏记、小录之书，皆记即日当时之事，求诸国史，最为实录。然皆言多\zd{鄙}朴，事罕圆备，\uline{终不能成其不刊，徒为后生作者削稿之资焉。}逸事者，皆前史所遗，后人所记，求诸异说，为益实多。及妄者为之，则苟载传闻，而无铨择，由是真伪不别，是非相乱，此其为弊之甚者也。
\end{materialtext}
\begin{flushright}（节选自《史通》，有删改）\end{flushright}

[注]\xq{1}赵正：秦始皇嬴政。《赵正书》收录于《北京大学藏西汉竹书》，西汉早期文献，2015年第一次对外出版。\xq{2}《秦二世元年十月甲午诏书》：木牍，2013年出土于湖南益阳兔子山。

\question{10}{材料一画波浪线的部分有三处需要断句，请用铅笔将答题卡上相应位置的答案标号涂黑。（3分）}

\duanju{昔者A秦王赵正出B游C天下D还E至柏F人G而病H病笃}

\question{11}{下列对材料中加点的词语及相关内容的解说，不正确的一项是（3分）}
\begin{choices}
\item 遽，立刻，与《促织》中“遽扑之，入石穴中”的“遽”意义、用法均相同。
\item 明，彰显，与《大学之道》中“在明明德”第一个“明”意义、用法均相同。
\item 参乘，陪乘，在此作动词，与《鸿门宴》中“沛公之参乘樊哙”的“参乘”用法不同。
\item 鄙，粗浅、浅陋，与《烛之武退秦师》中“越国以鄙远”的“鄙”意义不同。
\end{choices}

\question{12}{下列对材料有关内容的概述，不正确的一项是（3分）}
\begin{choices}
\item 材料一，赵正临终以“牛马”“蚊虻”设喻，指出大臣纷争会致两败俱伤，自己对朝局和百姓非常担忧，因而询问李斯对立太子的意见。
\item 材料二，诏书文告历数了秦二世即位以来的各项事务，并告诫官员不要用徭役赋税烦扰百姓，也不要因琐碎小事苛刻地弹劾县吏。
\item 材料三，赵高、胡亥、李斯拆毁了秦始皇所立的诏书，用假诏帮助胡亥成为太子，还列举公子扶苏和蒙恬的罪状，赐死了他们。
\item 材料四，逸事可以补充前代史书遗漏未载的内容，但若对传闻不加以甄别选择，是非真伪便会混淆，也会产生更为严重的危害。
\end{choices}

\question{13}{把材料中画横线的句子翻译成现代汉语。（8分）}

（1）元年与黔首更始，尽为解除故罪。

（2）终不能成其不刊，徒为后生作者削稿之资焉。

\question{14}{材料一与材料三对秦始皇临终传位的记载有哪些对立的信息？（4分）}

\par\vspace{2.2em}

\subsechead{阅读Ⅳ（9分）}
阅读下面这首宋诗，完成15～16题。

\begin{center}{\bfseries 垂丝钓}\par\vspace{0.15em}南宋·丘密\end{center}

\begin{poem}
戊戌\xq{1}近客，自入淮南，多所感怆作。\\
夕烽戍鼓，悲凉江岸淮浦。雾隐孤城，水荒沙聚。人共语，尽向来胜处，谩怀古。\\
问柳津花渡，露桥夜月，吹箫人在何许？缭墙禁籞\xq{2}，粉黛成黄土，惟有江东注。都无虏\xq{3}，似旧时得否？
\end{poem}

[注]\xq{1}淳熙五年戊戌（1178年），词人来到战后十四年的扬州城，怆然有感而作此词。\xq{2}禁籞（yù）：帝王的禁苑。\xq{3}都无虏：此句意为没有了金兵。

\question{15}{下列对这首词的理解和赏析，不正确的一项是（3分）}
\begin{choices}
\item 上阕以“夕烽”“戍鼓”开篇，点明烽烟依旧的时代背景，渲染了悲凉的气氛，为全词奠定了感情基调。
\item “雾隐孤城，水荒沙聚”运用白描手法，勾勒出孤城隐于暮霭的破败景象，令人想见扬州兵祸之深重。
\item 词人伫立江岸，极目远望柳津花渡、露桥夜月之下，吹箫人隐约可见，这与“夕烽戍鼓”的现实形成反差。
\item 下阕“缭墙禁籞”三句，与《扬州慢》中“念桥边红药，年年知为谁生”相似，以自然永恒反衬人事沧桑。
\end{choices}

\question{16}{“谩”与“问”二字贯通全词，这在情感表达上有何作用？请简析。（6分）}

\par\vspace{2.2em}

\subsechead{（五）名篇名句默写（本题共1小题，6分）}

\question{17}{补写出下列句子中的空缺部分。（6分）}

（1）“夜月”是古诗词中的常见意象。李白《蜀道难》中“\blank[6.5em]，\blank[6.5em]”两句，为了烘托蜀道的凄凉、表达羁旅之愁，也用到了这个词语。

（2）边塞诗中常将士兵和思妇对举，如《燕歌行》中“\blank[6.5em]，\blank[6.5em]”。

（3）在古文中，也有以问句结尾，引发读者思考的例子，比如《伶官传序》中“夫祸患常积于忽微，\blank[6.5em]，\blank[6.5em]？”

\sechead{二、语言文字运用（17分）}
阅读下面的文字，完成18～21题。

{\ywkai AI本身的局限，体现在中文学习的几个核心环节，构成从语用、文化到认知的多重短板。语用层面，汉语重“意合”，语义深藏于场景之中。例如，“坏了坏了，这下好了”与“好了好了，这下坏了”两句，前者常见于\underline{\hspace{0.6em}A\hspace{0.6em}}的小疏漏，系“先承认坏的结果，再故作轻松地自我安慰”；后者则是幻想破灭、大祸将至的麻烦，是“先抱有希望，再\underline{\hspace{0.6em}B\hspace{0.6em}}”。然而，AI仅凭词汇概率演算，只能识别两句均含负面情绪，却无法捕捉说话人心理落差的微妙差异，常将两者简单判定为同等程度的“糟心事”。\wave{\xq{1}在中文学习场景中，\xq{2}学习者若依赖AI来解读语句情绪，\xq{3}便很难把握精准的日常表达里细腻的分寸感，\xq{4}长期下来，\xq{5}将导致中文语感的缺失甚至损害。}在文化层面，AI还会因训练数据偏见而“自信断言不存在的事实”，可能导致学习者对中国文化产生错误认知。一个无法辨别“林黛玉倒拔垂杨柳”之荒谬的学习者，无论AI辅助下的口语多么流利，都不能被视为真正掌握了中文。\par}

\question{18}{汉语重“意合”，语义深藏于场景之中。文学作品中也不乏其例，下面的相关解说不正确的一项是（3分）}
\begin{choices}
\item “你是萍，——凭，——凭什么打我的儿子？”——“萍”“凭”二字的谐音转换，可见鲁侍萍悲愤复杂的心理，也说明她的理智战胜了感情。
\item “况且这通身的气派，竟不像老祖宗的外孙女儿，竟是个嫡亲的孙女……”——王熙凤见面就夸黛玉，模糊褒贬对象，表现出“一箭三雕”的泼辣。
\item “你放着罢，祥林嫂！”——这句话不符合正常语序，流露了四婶的慌乱和严厉，意在阻止祥林嫂参与年终祭祀，也是对她致命的判决。
\item “你总是很积极的。”——预感要离别，水生嫂对丈夫的支持、不舍、担忧与嗔怪交织在这句看似平淡的话中，含蓄而意蕴丰富。
\end{choices}

\question{19}{请在文中画横线处填入恰当的成语。（4分）}

\par\vspace{2.2em}

\question{20}{文中画波浪线的部分有两处表述不当，请指出其序号并修改，使语言准确流畅，逻辑严密，可少量增删词语，不得改变原意。（4分）}

\par\vspace{2.2em}

\question{21}{AI可能难以理解“林黛玉倒拔垂杨柳”这一热梗的趣味，请结合材料简析原因。（6分）}

\par\vspace{2.2em}

\sechead{三、写作（60分）}

\question{22}{阅读下面的材料，根据要求写作（60分）}

{\ywkai 本卷阅读II中，作者说“我不知道露天电影消失的时代好不好。想着让我难受。”}

时代与生活不断前行，这样的“消逝”似乎也难以避免。请你以一名青年朋友的身份给作者写一封信，谈谈你的感想与思考，回应他的困惑。

要求：选准角度，确定立意，自拟标题，完成写作任务；不要套作，不得抄袭；统一以“你的朋友晓华”为写信人，不得泄露个人信息；不少于800字。

\end{document}
